% VI49-DETAILED-HINTS
% VI49-DETAILED-HINT-01
\paragraph{Exercise VI.49.1.}
\begin{hint}
Restriction of a flabby sheaf is flabby.
\end{hint}

% VI49-DETAILED-HINT-02
\paragraph{Exercise VI.49.2.}
\begin{hint}
Skyscraper sheaves are flabby.
\end{hint}

% VI49-DETAILED-HINT-03
\paragraph{Exercise VI.49.3.}
\begin{hint}
Use a common power and a unit-ideal relation.
\end{hint}

% VI49-DETAILED-HINT-04
\paragraph{Exercise VI.49.4.}
\begin{hint}
There are only degree zero and one terms.
\end{hint}

% VI49-DETAILED-HINT-05
\paragraph{Exercise VI.49.5.}
\begin{hint}
Localize the homology at both covering elements.
\end{hint}

% VI49-DETAILED-HINT-06
\paragraph{Exercise VI.49.6.}
\begin{hint}
Apply affine vanishing.
\end{hint}

% VI49-DETAILED-HINT-07
\paragraph{Exercise VI.49.7.}
\begin{hint}
Use prime ideals and avoidance of a product.
\end{hint}

% VI49-DETAILED-HINT-08
\paragraph{Exercise VI.49.8.}
\begin{hint}
Finite intersections remain principal affine.
\end{hint}

% VI49-DETAILED-HINT-09
\paragraph{Exercise VI.49.9.}
\begin{hint}
Use the diagonal closed immersion.
\end{hint}

% VI49-DETAILED-HINT-10
\paragraph{Exercise VI.49.10.}
\begin{hint}
The Leray Čech complex has length two.
\end{hint}

% VI49-DETAILED-HINT-11
\paragraph{Exercise VI.49.11.}
\begin{hint}
Use the two-affine vanishing corollary.
\end{hint}

% VI49-DETAILED-HINT-12
\paragraph{Exercise VI.49.12.}
\begin{hint}
Use homogeneous quadrics.
\end{hint}

% VI49-DETAILED-HINT-13
\paragraph{Exercise VI.49.13.}
\begin{hint}
There are no compatible nonnegative exponents.
\end{hint}

% VI49-DETAILED-HINT-14
\paragraph{Exercise VI.49.14.}
\begin{hint}
Use the Laurent quotient.
\end{hint}

% VI49-DETAILED-HINT-15
\paragraph{Exercise VI.49.15.}
\begin{hint}
Keep the missing exponents \(n+1,\ldots,-1\).
\end{hint}

% VI49-DETAILED-HINT-16
\paragraph{Exercise VI.49.16.}
\begin{hint}
Use \(h^0=0\) and \(h^1=3\).
\end{hint}

% VI49-DETAILED-HINT-17
\paragraph{Exercise VI.49.17.}
\begin{hint}
Read the long exact sequence.
\end{hint}

% VI49-DETAILED-HINT-18
\paragraph{Exercise VI.49.18.}
\begin{hint}
Control \(H^1\) by surjectivity on global sections.
\end{hint}

% VI49-DETAILED-HINT-19
\paragraph{Exercise VI.49.19.}
\begin{hint}
Use four consecutive long-exact terms.
\end{hint}

% VI49-DETAILED-HINT-20
\paragraph{Exercise VI.49.20.}
\begin{hint}
Filter by skyscraper quotients.
\end{hint}

% VI49-DETAILED-HINT-21
\paragraph{Exercise VI.49.21.}
\begin{hint}
The point quotient is acyclic.
\end{hint}

% VI49-DETAILED-HINT-22
\paragraph{Exercise VI.49.22.}
\begin{hint}
The theorem requires quasi-coherence.
\end{hint}

% VI49-DETAILED-HINT-23
\paragraph{Exercise VI.49.23.}
\begin{hint}
Affine intersections must remain affine.
\end{hint}

% VI49-DETAILED-HINT-24
\paragraph{Exercise VI.49.24.}
\begin{hint}
Compare affine vanishing with asymptotic positive twisting.
\end{hint}

